\qhead{12. Open Question — The most important pricing topic for your company not covered above: geographic price discrimination (China vs.\ Europe).}

The pricing topic most central to BYD that the other answers do not fully capture is \textbf{third-degree price discrimination by geography} — selling comparable cars at very different prices across markets, fenced by location \parencite{nagle2018}.

\textbf{Evidence.} The Seagull sells for about \$7{,}800 in China \parencite{electrekSeagull}; in Europe the cheapest BYD (Dolphin Surf) starts near \euro19{,}990. The result is a \textbf{28.1\% overseas auto gross margin versus 17.2\% in China} in 2025 \parencite{byd2025}: Europe is the profit sanctuary, China the share battlefield.

\textbf{Why the fence holds.} Geography, vehicle homologation and the EU tariff together create a natural arbitrage barrier — a buyer cannot simply import a Chinese-spec Seagull into Germany — so the gap is sustainable without alienating either segment. Within each market the prices sit in legitimate \textbf{price bands} explained by tariffs, logistics, local cost-to-serve and the competitive reference price, which is ``expected'' band behaviour rather than ``outlaw'' pricing.

\textbf{Tariff and policy response.} The 17\% countervailing duty raises BYD's landed cost, but because European exchange value is far higher (set against Tesla and VW), BYD still earns a fat margin. Its strategic response is the new \textbf{Hungarian plant}, which localises production to escape the tariff, shorten logistics and protect the European price band \parencite{fortuneHungary}; selling increasingly through its own showrooms lets it hold ``fixed prices, flexible offers'' rather than discount away that margin.

\textbf{Significance.} Geographic discrimination is generally legal and often welfare-improving (it serves more buyers), limited only by Article~102 should BYD become dominant. I treat it as the most important topic because this margin gap is the \textbf{financial engine} that funds the domestic price war — the cross-subsidy without which the penetration strategy of the earlier answers would be unsustainable.
